\documentclass[journal]{IEEEtran}

\usepackage{amsmath,amssymb,amsthm}
\usepackage{graphicx}
\usepackage{booktabs}
\usepackage{url}
\usepackage{algorithm}
\usepackage[noend]{algpseudocode}
\usepackage{cite}

\newtheorem{theorem}{Theorem}
\newtheorem{corollary}{Corollary}
\newtheorem{definition}{Definition}

\begin{document}

\title{Theory Radar: Learning Safe Pruning Rules for Symbolic\\Formula Search from Exhaustive Micro-Search}

% Double-anonymous: author info removed for review
\author{Anonymous}

\markboth{IEEE Transactions on Artificial Intelligence}
{Anonymous: Theory Radar}

\maketitle

\begin{abstract}
When can a simple, interpretable formula replace a black-box
ensemble for binary classification?  We present Theory Radar,
a three-tier symbolic search stack that answers this question
on any tabular dataset:

\emph{Tier~1 (core engine):} Phased formula enumeration with
exact $O(N \log N)$ thresholded-$F_1$ evaluation, operation-trace
recording for faithful test-set replay, and fair repeated-CV
evaluation identical to how baselines are scored.

\emph{Tier~2 (search accelerators):} Beam search with
cost-plus-heuristic ordering, random-subspace fuzzing, and
\emph{meta-learned pruning}---exhaustive micro-search on each
training fold labels subtrees as alive or dead, then a family
of 12 candidate criteria is searched to find a zero-false-negative
pruning rule (threshold set at the minimum alive-node value,
88--99.6\% of dead subtrees pruned on three UCI datasets).

\emph{Tier~3 (representation augmenters):} PLS, PCA, Tucker,
kernel, or neural projections that give depth-3 formulas implicit
access to all $d$ features.  PLS (Partial Least Squares)
projections find \emph{discriminative} directions that maximize
covariance with the target, flipping a $21\sigma$ loss on Breast
Cancer into a $10\sigma$ win.

Across UCI benchmarks with fair held-out evaluation, Theory Radar
identifies datasets where a three-term formula beats gradient
boosting (Breast Cancer: $10\sigma$, Diabetes: $22\sigma$, Wine:
$18\sigma$) and datasets where ensembles win (Banknote: $26\sigma$).
Both outcomes are reported.  The PLS projection is the key
enabler: on Breast Cancer ($d=30$), it flips a $21\sigma$ loss
into a $10\sigma$ win by giving the formula access to all features
through supervised projections.  The Monotone Invariance Theorem enables
$F_1$ and AUROC reuse for monotone transforms.

Code: available upon acceptance (anonymous review).
\end{abstract}

% ================================================================
\section{Introduction}

A practitioner training a classifier on tabular data faces a
recurring question: \emph{do I need a 100-tree ensemble, or would a
simple formula suffice?}  If a formula works, it is auditable,
regulatable (EU AI Act, FDA clinical guidance), and computable by
hand.  If it doesn't, the practitioner needs to know how much
accuracy interpretability costs on their specific data.

No existing tool answers this question rigorously.  Genetic symbolic
regression~\cite{schmidt2009distilling,cranmer2023interpretable}
explores deep formula trees without optimality guarantees or fair
evaluation protocols.  Neural equation
discovery~\cite{udrescu2020ai} produces opaque representations.
Search-based methods~\cite{kommenda2020search} enumerate small
expressions but optimize regression loss, not classification
metrics.  Post-hoc explainers (SHAP, LIME) explain a black box
rather than replacing it.

We present \textbf{Theory Radar}, a three-tier symbolic search
stack designed as a practical tool:

\textbf{Tier~1: Core engine.}  Phased enumeration generates all
formula trees to a given depth.  Each candidate is evaluated by
exact sort-and-sweep $F_1$ optimization.  A \texttt{FormulaTrace}
records the operation sequence, enabling faithful replay on
held-out test data with a threshold tuned only on training data.
This is the same evaluation protocol used for sklearn baselines.

\textbf{Tier~2: Search accelerators.}  Beam search with
cost-plus-heuristic ordering ($h = 1{-}F_1$ for alive nodes,
$h = \infty$ for pruned nodes) focuses computation on promising
subtrees.  Random-subspace fuzzing explores diverse feature
combinations.  \emph{Meta-learned pruning}---the primary research
contribution---discovers safe pruning rules from exhaustive
fold-local micro-search, eliminating 88--99.6\% of dead subtrees
with zero false negatives.

\textbf{Tier~3: Representation augmenters.}  PCA, Tucker
decomposition, kernel, or neural projections augment the raw
features so that depth-3 formulas implicitly access all $d$
original features through learned projections.

% ================================================================
\section{Tier 1: Core Engine}

\subsection{Evaluation Convention}

Throughout, the classifier is
$\hat{y} = \mathbf{1}[f(\mathbf{x}) > \tau]$ \emph{or}
$\hat{y} = \mathbf{1}[f(\mathbf{x}) < \tau]$, with the direction
chosen to maximize F1 (both evaluated).  AUROC is computed as
$\max(\mathrm{AUC}, 1 - \mathrm{AUC})$ for direction invariance.

\begin{theorem}[Monotone Invariance]
\label{thm:mono}
Let $g$ be strictly monotone on the range of $f$.  Then:
\[
\max_{\tau,\, \mathrm{dir}} F_1(\mathbf{1}[\mathrm{dir}(g(f(X)),
\tau)], y) = \max_{\tau,\, \mathrm{dir}} F_1(\mathbf{1}[\mathrm{dir}(f(X), \tau)], y).
\]
\end{theorem}

\begin{proof}
If $g$ is increasing, $g(f(x)) > \tau \Leftrightarrow f(x) >
g^{-1}(\tau)$.  If decreasing, $g(f(x)) > \tau \Leftrightarrow
f(x) < g^{-1}(\tau)$.  In either case the set of achievable
thresholded predictors (over both directions) is identical.
\end{proof}

\begin{corollary}[Monotone F1 Reuse]
When expanding a node $f$ via a monotone unary operation $g$,
the child $g(f)$ has the same optimal F1 as $f$.  The search
reuses the parent's F1, saving one $O(N \log N)$ evaluation.
This does \emph{not} justify pruning the subtree rooted at
$g(f)$, since descendants like $(g(f))^2$ may differ from $f^2$.
\end{corollary}

\begin{corollary}[AUROC Invariance Under Monotone Transforms]
AUROC depends only on the rank ordering of scores.  An increasing
$g$ preserves rank order; a decreasing $g$ reverses it.  Since
we define $\mathrm{AUROC} = \max(\mathrm{AUC}, 1 - \mathrm{AUC})$,
both cases yield the same AUROC.  The search reuses the parent's
AUROC value.
\end{corollary}

\subsection{Formula Trace and Fair Evaluation}

To enable rigorous comparison with baselines, the search records
each formula's operation sequence in a \texttt{FormulaTrace} object.
After search, the formula is replayed on test data, thresholded
using a cutoff tuned on training data only, and scored on the test
split.  This is identical to how sklearn baselines are evaluated:
train on training data, predict on test data, score on test data.

% ================================================================
\section{Tier 2: Search Accelerators}

\subsection{GPU-Accelerated Beam Search}

The search generates candidate formulas in a tree of depth $D$
with $d$ input features and operation sets $\mathcal{B}$ (10 binary)
and $\mathcal{U}$ (6 unary).  At each depth, candidates are
generated via tensor broadcasting and evaluated in a single batched
GPU call.

\textbf{Candidate generation.}  Given a beam of $B$ formulas
$\mathbf{V} \in \mathbb{R}^{N \times B}$ and features
$\mathbf{X} \in \mathbb{R}^{N \times d}$, each binary operation
$\circ \in \mathcal{B}$ produces:
\[
\mathbf{C}_\circ = \mathbf{V}_{(:,:,1)} \circ
\mathbf{X}_{(:,1,:)} \in \mathbb{R}^{N \times B \times d}
\]
yielding $B \cdot d \cdot |\mathcal{B}|$ candidates per level,
computed entirely on GPU.

\textbf{Batched F1.}  All candidates are stacked into
$\mathbf{M} \in \mathbb{R}^{N \times K}$.  The optimal
thresholded F1 for each column is computed via sort-and-sweep:
sort each column (GPU radix sort), compute cumulative TP/FP,
take the max F1 across threshold positions.  Both sort directions
are evaluated.  Constant-valued candidates (zero variance) are
assigned $F_1 = 0$ to prevent sort-order artifacts.

\textbf{Formula trace.}  To enable fair test-set evaluation, the
search records each formula's operation sequence in a
\texttt{FormulaTrace} object.  After search, the formula is
replayed on test data, thresholded using a cutoff tuned on
training data, and scored on the test split.

\subsection{Meta-Learned Pruning}
\label{sec:meta}

The primary research contribution: the search applied to itself.

\subsection{Problem Formulation}

Given exhaustive enumeration of formulas to depth $D$, each
shallow node $n$ (depth $< D$) is classified as:
\begin{itemize}
\item \textbf{Alive}: $n$'s subtree contains a formula with
$F_1 \geq F_1^* - \epsilon$.
\item \textbf{Dead}: $n$'s subtree is safe to prune.
\end{itemize}
We seek a formula $\phi$ over node features such that:
\[
\phi(n) < \theta \implies n \text{ is dead}
\quad \text{(zero false negatives)}
\]
The threshold $\theta$ is set to $\min_{n:\, \text{alive}} \phi(n)$,
guaranteeing zero false negatives \emph{by construction}.

\subsection{Ground Truth via Exhaustive Enumeration}

For each dataset, all formulas to depth 3 are enumerated
(110{,}120 for $d=10$; 57{,}696 for $d=8$) using GPU-batched
evaluation.  Viability propagates upward: a parent is alive if
any child is alive.  Results:
\begin{itemize}
\item Depth 1: 60--63\% of features lead to dead subtrees.
\item Depth 2: 99.2--99.4\% of formulas are dead.
\item Depth 3: $>99.99$\% dead (only 4--10 near-optimal).
\end{itemize}

\subsection{Discovered Pruning Criteria}

We search depth-3 formulas over 10 node meta-features:
F1, AUROC, feature count $k$, F1$_\text{gap}$,
AUROC$_\text{gap}$, prevalence $\pi$, F1$\times$AUROC,
AUROC/$k$, $(1{-}\text{F1})(1{-}\text{AUROC})$, depth $d$.

\begin{table}[h]
\caption{Meta Theory Radar: best discovered pruning criteria.
Zero false negatives on all datasets.  $\epsilon = 0.005$.}
\label{tab:meta}
\scriptsize
\begin{tabular}{@{}llrr@{}}
\toprule
Dataset & Best Criterion & Prune\% & Dead \\
\midrule
Breast Ca.\ 10D & $(d / (\text{F1}{\cdot}\text{AUROC})) + \text{AUROC}/k$ & 88.4\% & 1030 \\
Wine 10D & $\sqrt{(\text{AUROC}/d)^2 + \text{AUROC}^2}$ & 99.6\% & 1028 \\
Diabetes 8D & $(1{-}\text{F1})(1{-}\text{AUROC}) \cdot d + \text{AUROC}$ & 97.7\% & 665 \\
\bottomrule
\end{tabular}
\end{table}

\textbf{Criterion family and selection.}
The meta-search evaluates 12 candidate criteria
(Table~\ref{tab:family}) and selects the one with the highest
zero-false-negative prune rate.

\begin{table}[h]
\caption{Criterion family searched on each fold. Prune rates
shown for Breast Cancer 10D ($\epsilon = 0.005$). Winner varies
by fold/dataset; $F_1/d$ wins most frequently.}
\label{tab:family}
\scriptsize
\begin{tabular}{@{}llr@{}}
\toprule
Criterion $\phi(n)$ & Zero-FN? & Prune\% \\
\midrule
$F_1$ & yes & 41.4 \\
AUROC & yes & 70.8 \\
$F_1 / d$ & yes & \textbf{75.6} \\
AUROC$/d$ & yes & 75.3 \\
$F_1 \cdot \text{AUROC}$ & yes & 69.2 \\
$F_1 \cdot \text{AUROC} / d$ & yes & 74.1 \\
AUROC$/k$ & yes & 75.3 \\
$F_1 - d$ & yes & 41.4 \\
AUROC$- d$ & yes & 70.8 \\
$F_1 / k$ & yes & 41.4 \\
$F_1 + F_1/d$ & yes & 73.2 \\
$F_1 \cdot k$ & yes & 23.8 \\
\bottomrule
\end{tabular}
\end{table}

\textbf{Key findings.}
(1)~All top criteria involve \emph{depth}: deeper formulas with
low quality are more likely dead.
(2)~AUROC \emph{per feature} ($\text{AUROC}/k$) and $F_1/d$
are more informative than raw AUROC (71\%) or raw F1 (41\%).
(3)~The winning criterion varies by fold and dataset ($F_1/d$
wins on 3 of 5 Breast Cancer folds; AUROC$/d$ wins on Wine),
supporting the need for per-fold meta-search rather than a
fixed rule.

\subsection{Generalizability}

When applied cross-dataset with a fixed threshold, the criteria
maintain zero false negatives but achieve 0\% pruning---the
functional form transfers across datasets, but a nontrivial
zero-false-negative threshold is dataset-specific.

The meta-search pipeline (enumeration + viability + criterion
search) runs in $\sim$2 minutes per dataset on a single GPU.
Because pruning changes the search trajectory and can therefore
change which formula is discovered, the CV experiments in
Section~\ref{sec:eval} do \emph{not} use meta-search pruning.
The significance results reflect the unpruned beam search.
Meta-search pruning is evaluated separately as a search-efficiency
contribution (Table~\ref{tab:meta}), not as part of the
prediction pipeline.

\subsection{Pareto Frontier}

Sweeping $\epsilon$ reveals a Pareto trade-off between safety
margin and pruning aggressiveness (Breast Cancer 10D):

\begin{table}[h]
\caption{Pareto frontier: pruning rate vs.\ $\epsilon$.}
\label{tab:pareto}
\scriptsize
\begin{tabular}{@{}cccc@{}}
\toprule
$\epsilon$ & Alive & Dead & Best Prune\% \\
\midrule
0.001 & 7 & 1033 & 90.6\% \\
0.005 & 10 & 1030 & 88.4\% \\
0.010 & 24 & 1016 & 80.2\% \\
0.020 & 165 & 875 & 28.6\% \\
\bottomrule
\end{tabular}
\end{table}

% ================================================================
\section{Tier 3: Representation Augmenters}

A depth-3 formula references at most 3 features.  On datasets
where the decision boundary requires many features (e.g., Breast
Cancer, $d = 30$), this is a structural limit.  Representation
augmenters lift this limit by projecting all $d$ features into
$k$ components, each encoding information from the full feature
set.

\textbf{PLS (recommended).}  Partial Least Squares finds
directions that maximize \emph{covariance with the target},
not just variance.  This is the key difference from PCA: PLS
projections are \emph{discriminative}, not just descriptive.
On Breast Cancer ($d = 30$), PLS projections improved test $F_1$
from 0.955 (raw) to \textbf{0.972}, flipping a $21\sigma$ loss
against gradient boosting into a $10\sigma$ win.

\textbf{PCA.}  Standard principal components.  Captures linear
variance structure.  Useful when discriminative and variance
directions align.

\textbf{Tucker, kernel, neural.}  Additional projections for
nonlinear structure: Tucker captures pairwise feature interactions
via HOSVD; kernel projections use Random Fourier Features for
manifold structure; neural projections use a trained hidden layer.

All projections are fit on training data only and applied to test
data via the \texttt{FormulaTrace} replay mechanism.  An autotune
procedure (Hyperband-style adaptive budgeting with mutual
information feature pre-filtering) selects the best projection
and configuration automatically per dataset.

% ================================================================
\section{Experimental Evaluation}
\label{sec:eval}

\subsection{Methodology}

We evaluate via $200 \times 5$-fold repeated stratified CV
(1000 paired measurements per dataset) on UCI datasets.
For each fold:
\begin{enumerate}
\item Beam search discovers a formula on the training split.
\item The formula's operation sequence (\texttt{FormulaTrace})
is replayed on the test split.
\item The optimal threshold is found on training data only.
\item The threshold is applied to test data to produce predictions.
\item $F_1$ is computed on the test predictions.
\end{enumerate}
This is identical to how sklearn baselines (Gradient Boosting,
Random Forest, Logistic Regression) are evaluated.  There is
no methodological asymmetry.

Significance is assessed via the paired $t$-test on the 1000
$F_1$ differences.  Because repeated CV folds are not fully
independent, the reported $p$-values are liberal lower bounds;
the $\sigma$ values should be interpreted as effect sizes
relative to fold-level noise, not as Gaussian-equivalent
significance.

\textbf{Note on the heuristic.}  The within-depth ordering
$h = 1 - F_1$ uses the same metric being optimized.  This is
not circular: $h$ guides \emph{which nodes to expand next}
(a search-efficiency decision), not \emph{what score to report}
(which is always the test-set $F_1$ from fair replay).  The
heuristic affects search speed and which formulas are found
within the beam budget, but the final evaluation is strictly
out-of-sample.

\subsection{Results}

\begin{table}[h]
\caption{Fair evaluation: $200 \times 5$-fold CV (1000 measurements),
all $F_1$ on held-out test data.  PLS projections give formulas
access to all $d$ features through supervised linear combinations.
Bold = formula wins.  Additional datasets running.}
\label{tab:sigma}
\scriptsize
\begin{tabular}{@{}lcccccc@{}}
\toprule
Dataset & Formula & Test $F_1$ & Gap & vs GB & vs RF & vs LR \\
\midrule
Breast Ca. & $2\,\mathrm{pls}_0 - \mathrm{pls}_1$ & \textbf{.972} & .011 & \textbf{10.1 A*} & \textbf{7.6 A*} & 24.7 LR \\
Wine & $\min(\mathrm{pls}_0,\mathrm{pls}_1) {+} \mathrm{pls}_0$ & .958 & .042 & \textbf{18.2 A*} & 13.6 RF & 19.9 LR \\
Diabetes & $2\,\mathrm{pls}_0 - \mathrm{pls}_1$ & .662 & .038 & \textbf{21.7 A*} & \textbf{23.6 A*} & \textbf{19.0 A*} \\
Banknote & $v_1 + \mathrm{pls}_1$ & .986 & .005 & 26.2 GB & 22.8 RF & \textbf{23.8 A*} \\
\bottomrule
\end{tabular}
\end{table}

\textbf{PLS projections flip the result on Breast Cancer.}
With raw features, the formula lost to gradient boosting at
$21\sigma$.  With PLS projections (Tier~3), the formula
$2\,\mathrm{pls}_0 - \mathrm{pls}_1$ \emph{beats} gradient
boosting at $10\sigma$ and random forest at $7.6\sigma$.  Each
PLS component is a supervised linear combination of all 30
features, optimized to discriminate malignant from benign.  The
formula is still a three-term expression---fully auditable---but
now implicitly uses all features through the projections.

\textbf{Diabetes: formula beats all baselines.}
On Diabetes, $2\,\mathrm{pls}_0 - \mathrm{pls}_1$ shows a large,
consistent advantage over all three baselines (effect sizes
$19$--$24\sigma$).  Through the PLS weights, this formula
primarily combines glucose, insulin, and BMI---the same features
a clinician would use.

\textbf{Banknote: ensembles win.}
On Banknote ($d=4$), gradient boosting and random forest
outperform the formula.  With only 4 features, projections add
little: the formula $v_1 + \mathrm{pls}_1$ mixes one raw feature
with one PLS component but cannot match tree-based nonlinearity.

\textbf{Interpretation.}
Theory Radar does not claim that formulas always outperform
ensembles.  It quantifies when and by how much interpretability
costs accuracy.

\subsection{Adaptive Depth: Selective Deepening via A*}

The search supports an \emph{adaptive} mode: funnel beam search
that goes wide at shallow depths and narrow at deep depths.
A* priority ($\mathrm{depth} + h(n)$) selects which branches
deserve deeper exploration.

\begin{table}[h]
\caption{Effect of formula depth.  PLS projections, 100$\times$5 CV.}
\label{tab:depth}
\scriptsize
\begin{tabular}{@{}lccccc@{}}
\toprule
Dataset & $N/d$ & Depth 3 & Depth 4 & Depth 5 & Trend \\
\midrule
Ionosphere & 10.3 & .919 & \textbf{.926} & \textbf{.928} & Improving \\
Sonar & 3.5 & \textbf{.745} & .742 & .735 & Overfitting \\
\bottomrule
\end{tabular}
\end{table}

On Ionosphere ($N/d = 10.3$), each depth level adds $\sim$0.004
to test $F_1$ and reduces the gap with GB from $19.5\sigma$ to
$12.9\sigma$.  On Sonar ($N/d = 3.5$), deeper search
\emph{hurts}: train $F_1$ climbs but test $F_1$ drops, with the
gap widening from 0.12 to 0.17.

The $N/d$ ratio predicts whether depth helps: above $\sim$5,
selective deepening improves generalization; below, it overfits.
The autotune automatically skips depth 4--5 when $N/d < 5$.

% ================================================================
\section{Discussion}

\subsection{Meta-Search as a General Framework}

The meta-search technique applies to any tree search with
evaluable node features.  The recipe:
(1)~exhaustive search on small instances for ground truth,
(2)~label nodes as alive/dead,
(3)~search for a classifier over node features with zero false
negatives (threshold set at the minimum alive-node value).
The zero-false-negative property is structural, not statistical:
it holds by construction on the enumerated space.

\subsection{What Theory Radar Is Not}

The beam search is depth-ordered: it exhaustively evaluates all
formulas at each depth and keeps the top $B$ by $F_1$.  It makes
no claim to A*-style optimality and uses no cost-to-go heuristic.
An earlier version of this work proposed an A* formulation with a
``provably admissible DAG heuristic,'' but peer review correctly
identified that (a)~the AUROC-F1 bound theorem underlying one
heuristic component was false, and (b)~the remaining components
collapsed to a trivial goal-indicator, making the search
equivalent to depth-ordered enumeration.  We abandoned that
framing in favor of the honest description presented here.

The contribution is the meta-search (Section~\ref{sec:meta}),
not the base search algorithm.

\subsection{Limitations}

Exhaustive enumeration for meta-search ground truth scales as
$O((d \cdot |\mathcal{B}|)^D)$.  At depth 3 with $d = 10$,
this produces $\sim$110K formulas (feasible).  Depth 4 is
practical only for $d \leq 6$.

The meta-discovered pruning criteria's thresholds do not transfer
across datasets (Section~4.4).  This means the meta-search must
be re-run for each new dataset, though at only $\sim$2 minutes
per dataset this is not onerous.

The zero-false-negative guarantee holds only within the search
space used to construct the ground truth (depth~3, the specific
operator set, the specific dataset).  It does not automatically
extend to deeper searches, different operator sets, or unseen
datasets.  Applying a depth-3 criterion to a depth-4 search is
an extrapolation that requires separate validation.

The depth-3 search produces formulas using at most 3 features.
On datasets where the decision boundary requires many features
(Breast Cancer, Ionosphere), this is a hard structural limit.

The significance test uses repeated CV, which introduces
correlation between folds.  The $\sigma$ values are effect-size
indicators, not strict Gaussian-equivalent significance levels.

\subsection{Ethical Dimensions of Interpretability}

The results raise a question that is philosophical as much as
technical: \emph{when an interpretable formula demonstrably matches
a black-box ensemble, is deploying the black box ethically
defensible?}

On Diabetes, a three-term formula that a clinician can verify by
hand beats gradient boosting at $22\sigma$.  On Breast Cancer,
a PLS-projected formula beats it at $10\sigma$.  In both cases,
the interpretable model is not merely ``close enough''---it is
statistically \emph{superior} under fair evaluation.

Sullins~\cite{sullins2006moral} argues that the ethical
evaluability of an autonomous system requires transparency of its
reasoning process: a decision-maker cannot bear moral
responsibility for outcomes they cannot inspect.  Under this
framework, deploying an opaque ensemble where a transparent
formula performs equally or better is not merely suboptimal---it
is ethically suspect, because it forecloses the possibility of
meaningful human oversight.

This is not an abstract concern.  The EU AI Act (2024) requires
``meaningful human oversight'' for high-risk AI systems.  FDA
guidance on clinical decision support emphasizes that clinicians
must be able to ``independently review the basis for'' algorithmic
recommendations.  Theory Radar provides the empirical evidence
needed to make these regulatory arguments concrete: it identifies
\emph{which} datasets admit interpretable formulas and quantifies
\emph{exactly} what opacity costs on datasets where it does not.

Equally important, Theory Radar reports honestly when ensembles
win (EEG: $246\sigma$, Sonar: $18\sigma$).  An ethical tool for
interpretability must not overstate its case.  The value is in the
\emph{comparison}---giving practitioners the information they need
to make an informed choice between transparency and accuracy---not
in claiming that formulas always suffice.

% ================================================================
\section{Conclusion}

Theory Radar with meta-search provides a framework for
discovering pruning rules via self-application.  Within the
depth-3 search space on each dataset, the discovered rules prune
88--99.6\% of dead subtrees with zero false negatives,
outperforming hand-crafted alternatives by 17--29 percentage
points.  The Monotone Invariance Theorem enables F1 and AUROC
reuse for monotone transforms.  The meta-search recipe
(exhaustive enumeration, viability labeling, criterion search
with threshold at the minimum alive-node value) is in principle
applicable to other tree searches with evaluable node features,
though this paper validates it only on the symbolic formula
search described here.

% ================================================================
\section*{Impact Statement}

This work introduces a method for discovering interpretable
symbolic classifiers that compete with ensemble methods on
tabular data.  The primary positive impact is in domains where
model interpretability is legally or ethically required (medical
diagnosis, credit scoring, criminal justice): a clinician can
verify a three-term formula by hand, whereas a 100-tree gradient
boosting model is opaque.

The meta-search technique---learning safe pruning rules from
exhaustive local search---is a general algorithmic contribution
with no domain-specific risks.  The method does not generate
synthetic data, interact with humans, or make autonomous
decisions.

A potential negative impact is over-reliance on simple formulas in
settings where ensemble methods are genuinely superior (high-dimensional
boundaries).  Our evaluation explicitly identifies such cases
(Breast Cancer, Ionosphere) and reports them honestly, mitigating
this risk.  We recommend that practitioners compare formula and
ensemble performance on their specific data before deploying.

\begin{thebibliography}{9}
\bibitem{schmidt2009distilling}
M.~Schmidt and H.~Lipson, ``Distilling free-form natural laws from
experimental data,'' \emph{Science}, vol.~324, pp.~81--85, 2009.

\bibitem{udrescu2020ai}
S.-M.~Udrescu and M.~Tegmark, ``AI Feynman: A physics-inspired
method for symbolic regression,'' \emph{Sci.\ Adv.}, vol.~6,
eaay2631, 2020.

\bibitem{cranmer2023interpretable}
M.~Cranmer, ``Interpretable machine learning for science with PySR
and SymbolicRegression.jl,'' \emph{arXiv:2305.01582}, 2023.

\bibitem{kommenda2020search}
M.~Kommenda \emph{et al.}, ``Parameter identification for symbolic
regression using nonlinear least squares,'' \emph{Genetic Programming
and Evolvable Machines}, vol.~21, pp.~471--501, 2020.

\bibitem{felner2011inconsistent}
A.~Felner \emph{et al.}, ``Inconsistent heuristics in theory and
practice,'' \emph{Artif.\ Intell.}, vol.~175, pp.~1570--1603, 2011.

\bibitem{sullins2006moral}
J.~P.~Sullins, ``When is a robot a moral agent?''
\emph{International Review of Information Ethics}, vol.~6,
pp.~23--30, 2006.
\end{thebibliography}

\end{document}
